\documentclass[12pt,a4paper]{article}

\usepackage[utf8]{inputenc}
\usepackage[T1]{fontenc}
\usepackage{lmodern}
\usepackage[margin=1in]{geometry}
\usepackage{amsmath,amssymb,amsfonts}
\usepackage{graphicx}
\usepackage{booktabs}
\usepackage{longtable}
\usepackage{array}
\usepackage{tabularx}
\usepackage{multirow}
\usepackage{hyperref}
\usepackage{cleveref}
\usepackage{natbib}
\usepackage{setspace}
\usepackage{enumitem}
\usepackage{xcolor}
\usepackage{float}
\usepackage{listings}
\usepackage{pdflscape}

\hypersetup{
    colorlinks=true,
    linkcolor=blue!70!black,
    citecolor=green!50!black,
    urlcolor=blue!70!black
}

\onehalfspacing

\title{A Cross-Jurisdiction Crosswalk of Machine-Learning Model-Risk Regulation, with an Open Compliance Benchmark for Quantitative Finance}

\author{
    Akram Khan\thanks{Independent Researcher. ORCID: \href{https://orcid.org/0009-0002-7521-8648}{0009-0002-7521-8648}. Email: \href{mailto:1819ak@gmail.com}{1819ak@gmail.com}. The crosswalk, the scoring harness, and the benchmark leaderboard are released at \href{https://github.com/ayk5511/mrm-crosswalk}{github.com/ayk5511/mrm-crosswalk}.}
}

\date{April 2026}

\begin{document}

\maketitle

\begin{abstract}
Machine-learning model-risk regulation in finance has fragmented across jurisdictions. The U.S.\ SR 11-7 / OCC Bulletin 2011-12, the EU AI Act (Regulation 2024/1689), the U.K.\ Prudential Regulation Authority's SS1/23, Singapore's MAS FEAT principles and Veritas assessment methodologies, and Canada's OSFI Guideline E-23 each impose overlapping but non-identical obligations on model documentation, data provenance, independent validation, per-prediction logging, record retention, monitoring, and governance. A firm deploying one model into several markets must currently reconcile these texts by hand, and a researcher auditing the published literature has no common yardstick. This paper makes two paired contributions. First, we construct a \emph{cross-jurisdiction crosswalk} that maps twelve audit-trail and model-risk-management requirements onto the specific provision of each regime (and onto the voluntary NIST AI RMF and ISO/IEC 42001 frameworks), making explicit where the regimes converge, where they diverge, and where a requirement in one has no analogue in another. Second, we release an \emph{open compliance benchmark} that scores a machine-learning pipeline against the union of these requirements, with a reproducible scoring harness and a public leaderboard. Both artefacts are released under permissive licences and versioned against the regulations' effective dates, so that practitioners, auditors, and future researchers can extend them as the law evolves.

\medskip
\noindent\textbf{Keywords:} model risk management, SR 11-7, EU AI Act, PRA SS1/23, OSFI E-23, MAS Veritas, compliance benchmark, audit trail, regulatory technology, machine learning

\medskip
\noindent\textbf{JEL Classification:} C88, G18, G28, K23
\end{abstract}

\newpage
\tableofcontents
\newpage

\section{Introduction}\label{sec:intro}

A bank that deploys a single credit-scoring or forecasting model into more than one market is now subject to several distinct model-risk regimes at once. In the United States, model risk management has been governed since 2011 by the joint Federal Reserve / OCC guidance SR 11-7 \citep{FedSR117, OCC2011}. In the European Union, the AI Act \citep{EUAIAct} imposes record-keeping, logging, and retention duties on high-risk AI systems, with the high-risk provisions taking effect on 2 August 2026. The United Kingdom's Prudential Regulation Authority issued its first dedicated model-risk supervisory statement, SS1/23 \citep{PRASS123}, effective 17 May 2024. Singapore's Monetary Authority published the FEAT principles \citep{MASFEAT} and the Veritas assessment methodologies \citep{MASVeritas}. Canada's OSFI finalised Guideline E-23 \citep{OSFIE23} in September 2025, effective 1 May 2027, expanding its scope to all models including machine learning. Two voluntary frameworks (the NIST AI Risk Management Framework \citep{NIST_AI_100_1} and ISO/IEC 42001 \citep{ISO42001}) sit alongside the binding regimes.

These regimes overlap heavily in intent but differ in letter. All require some form of documentation, validation, and governance; only the EU AI Act fixes a numerical record-retention floor; only OSFI E-23 makes model \emph{decommissioning} an explicit lifecycle stage; fairness assessment is central to the EU and Singapore regimes but absent from the original SR 11-7. A compliance officer reconciling these texts, and a researcher attempting to measure whether published work would survive any of them, both face the same obstacle: there is no common yardstick. The fragmentation is not academic: machine learning is now a core methodology in quantitative finance \citep{GuKellyXiu2020}, and the same model is routinely deployed across jurisdictions.

This paper supplies that yardstick and makes two paired contributions. First (\Cref{sec:crosswalk}), we construct a \emph{cross-jurisdiction crosswalk}: a single table mapping twelve audit-trail and model-risk-management requirements onto the specific provision of each of the five regulatory regimes and the two voluntary frameworks, so that the table can be read down a column for a single regime's full set of obligations or across a row to compare the regimes on one requirement. Second (\Cref{sec:benchmark}), we release an \emph{open compliance benchmark} that scores a machine-learning pipeline against the union of these requirements, with a reproducible scoring harness and a public leaderboard, so that compliance against the combined standard becomes a measurable, comparable quantity rather than a manual reading exercise.

\subsection{Contributions}\label{sec:contributions}

\begin{itemize}[leftmargin=*]
    \item A twelve-requirement crosswalk of machine-learning model-risk regulation across five regulatory regimes and two voluntary frameworks (\Cref{tab:crosswalk}), with each cell tied to the citable provision and a marking of whether each requirement is present, implied, partial, or absent in each regime.
    \item A structured analysis of cross-regime convergence and divergence: which requirements are universal, which are jurisdiction-specific, and which create genuine multi-market compliance conflicts.
    \item An open, versioned compliance benchmark and scoring harness, released under permissive licences, that scores an ML pipeline against the union standard and against any single regime, with a public leaderboard seeded on public-data pipelines.
\end{itemize}

\subsection{Scope and limitations}\label{sec:scope}

The crosswalk is an \emph{interpretive} mapping of public regulatory text, not legal advice; firms remain responsible for their own compliance determinations with their supervisors. Regulations evolve: the U.S.\ replaced SR 11-7 / OCC 2011-12 with new interagency guidance in 2026, and OSFI E-23 only takes effect in 2027. The crosswalk is therefore versioned to the texts as of the dates recorded in \Cref{sec:crosswalk}, and the released artefact is designed to be updated as the law changes.

\subsection{Roadmap}\label{sec:roadmap}

\Cref{sec:crosswalk} describes the crosswalk methodology, presents the mapping, and analyses convergence and divergence. \Cref{sec:benchmark} describes the open compliance benchmark, its scoring harness, and the seed leaderboard. \Cref{sec:discussion} discusses use, limitations, and future maintenance.

\section{The cross-jurisdiction crosswalk}\label{sec:crosswalk}

\subsection{Method}\label{sec:crosswalk_method}

We fix twelve requirements that recur across model-risk and AI-governance regimes and that together define an auditable model lifecycle: (1) model inventory and identification; (2) documentation sufficient for independent replication; (3) data provenance and lineage; (4) versioning and change control; (5) independent validation; (6) per-prediction or event-level logging; (7) record retention with a defined duration; (8) ongoing monitoring and drift detection; (9) governance and named accountability; (10) transparency and explainability; (11) fairness and bias assessment; and (12) decommissioning as an explicit lifecycle stage. These twelve extend the audit-trail rubric of \citet{KhanMrAudit2026} from a single-regime (SR 11-7 / EU AI Act) scope to the multi-regime setting. Several of them also have an established methodological literature: structured model documentation via Model Cards \citep{MitchellModelCards2019}, the reproducibility-in-AI scholarship of \citet{GundersenKjensmo2018}, and risk-data-aggregation standards such as BCBS~239 \citep{BCBS239} for data lineage.

For each requirement we read the public text of five binding or semi-binding regimes --- SR 11-7 / OCC 2011-12 \citep{FedSR117, OCC2011}, the EU AI Act \citep{EUAIAct}, PRA SS1/23 \citep{PRASS123}, the MAS FEAT principles and Veritas methodologies \citep{MASFEAT, MASVeritas}, and OSFI Guideline E-23 \citep{OSFIE23} --- and two voluntary frameworks (NIST AI RMF \citep{NIST_AI_100_1} and ISO/IEC 42001 \citep{ISO42001}), and record the most on-point provision. \Cref{tab:crosswalk} is the result. Cells cite the governing section, article, or principle; \textit{partial} marks a provision that addresses the requirement only in part; \textit{impl.} marks a requirement that is implied by the regime's lifecycle language without a dedicated clause; and ``---'' marks the absence of an explicit analogue (which is not the same as prohibition).

Two honesty caveats govern the table. It is an \emph{interpretive} reading of public text, not legal advice. And it is \emph{versioned}: SR 11-7 / OCC 2011-12 reflects the 2011 guidance as historically in force (rescinded and replaced by revised interagency guidance in 2026); the EU AI Act reflects Regulation 2024/1689 with high-risk duties effective 2 August 2026; PRA SS1/23 is effective 17 May 2024; OSFI E-23 is the September 2025 final text effective 1 May 2027.

\begin{landscape}
\begin{table}[p]
\centering
\scriptsize
\setlength{\tabcolsep}{4pt}
\renewcommand{\arraystretch}{1.25}
\caption{Cross-jurisdiction crosswalk of twelve model-risk / audit-trail requirements. Cells cite the governing provision; \textit{partial} = partially addressed, \textit{impl.} = implied by lifecycle language, ``---'' = no explicit analogue. NIST/ISO column gives the relevant NIST AI RMF function / ISO/IEC 42001 clause.}\label{tab:crosswalk}
\begin{tabular}{@{}p{3.0cm}p{2.6cm}p{3.0cm}p{2.4cm}p{2.6cm}p{2.4cm}p{2.6cm}@{}}
\toprule
\textbf{Requirement} & \textbf{SR 11-7 (US)} & \textbf{EU AI Act} & \textbf{PRA SS1/23 (UK)} & \textbf{MAS FEAT / Veritas (SG)} & \textbf{OSFI E-23 (CA)} & \textbf{NIST / ISO} \\
\midrule
1. Model inventory / identification & \S III & Art.~11, Annex~IV; registration Art.~49 & Principle 1 & --- & Explicit enterprise inventory w/ metadata & MAP-1 / Annex~A \\
2. Replicable documentation & \S IV--V & Art.~11, Annex~IV & Principle 3 & Veritas Doc 3C & Lifecycle documentation & MEASURE / \S 7.5 \\
3. Data provenance / lineage & \S IV (data quality) & Art.~10 & Principle 3 & Veritas Doc 3A (data) & Model-design data component & MAP-3 / \S 7.4; cf.\ BCBS~239 \\
4. Versioning / change control & \S V & Art.~12; re-conformity & Principles 2--3 & --- & Lifecycle (deploy / monitor) & MANAGE / \S 8 \\
5. Independent validation & \S V (core) & Art.~17, Art.~43 & Principle 4 (explicit) & Accountability & Independent model review & MEASURE-2 / \S 9 \\
6. Per-prediction / event logging & \S V (outcomes; partial) & Art.~12; items in Art.~12(3) (biometric) & --- & \textit{partial} & Monitoring (\textit{partial}) & --- \\
7. Record retention (duration) & --- & \textbf{Art.~19 ($\geq$6 months)} & --- & --- & Through lifecycle + post-decommission & --- \\
8. Ongoing monitoring / drift & \S V & Art.~72 (post-market), Art.~9 & Principles 3, 5 & Accountability & Model monitoring & MANAGE-4 / \S 9.1 \\
9. Governance / accountability & \S VI & Art.~26, Art.~22 & Principle 2 & Accountability (FEAT) & Governance, scaled to risk & GOVERN / \S 5 \\
10. Transparency / explainability & \S IV & Art.~13; Art.~86 (individual) & Principle 3 & Transparency (FEAT; Doc 3C) & Explainability considerations & MAP-2 / \S 6 \\
11. Fairness / bias assessment & --- & Art.~10(2), Art.~15 & --- & \textbf{Fairness (FEAT; Doc 3A)} & \textit{partial} & MEASURE-2.11 / --- \\
12. Decommissioning / lifecycle end & \textit{impl.} & \textit{impl.} & Principle 3 & --- & \textbf{Explicit decommission stage} & MANAGE / \S 8 \\
\bottomrule
\end{tabular}
\end{table}
\end{landscape}

\subsection{Where the regimes converge}\label{sec:converge}

Five requirements are effectively universal. \emph{Independent validation} (row 5), \emph{documentation sufficient for an independent party to understand and reproduce the model} (row 2), \emph{governance with named accountability} (row 9), \emph{ongoing monitoring} (row 8), and a \emph{model inventory} (row 1) appear, in some form, in every binding regime except where a regime is principles-only. SR 11-7's three pillars (development, validation, governance), PRA SS1/23's five principles, and OSFI E-23's lifecycle all reduce to the same backbone. A firm that satisfies this backbone has satisfied the common core of every regime studied.

\subsection{Where the regimes diverge}\label{sec:diverge}

Three requirements are jurisdiction-specific and are the source of genuine multi-market compliance gaps:

\begin{itemize}[leftmargin=*]
    \item \textbf{Record retention with a fixed floor (row 7)} is unique to the EU AI Act: Article 19 sets a $\geq$6-month minimum on automatically generated logs. No other regime fixes a duration. A pipeline tuned to SR 11-7 alone will typically have no retention policy that satisfies Article 19.
    \item \textbf{Fairness / bias assessment (row 11)} is central to the EU AI Act (Article 10's data-governance bias examination) and to MAS FEAT (the Fairness principle and Veritas Document 3A), but is absent from the original SR 11-7 and PRA SS1/23, which are prudential rather than conduct instruments.
    \item \textbf{Decommissioning (row 12)} is an explicit lifecycle stage only in OSFI E-23; elsewhere it is at most implied. A firm with no documented model-retirement process is compliant in the US and UK but not in Canada.
\end{itemize}

Per-prediction logging (row 6) is the sharpest example of \emph{convergence in spirit, divergence in letter}: the EU AI Act mandates event logging (Article 12, with the enumerated items of Article 12(3) applying specifically to biometric systems of Annex III(1)(a)), SR 11-7 reaches only outcomes analysis, and PRA SS1/23 is silent. A pipeline built for one regime's logging expectation will under- or over-comply with another's. The benchmark of \Cref{sec:benchmark} scores against the \emph{union}, so that a single instrumented pipeline can be shown to clear the highest bar across all regimes simultaneously.

\section{An open compliance benchmark}\label{sec:benchmark}

The crosswalk turns five regulatory texts into one structured object. The benchmark turns that object into a \emph{measurement}: given a machine-learning pipeline (a repository, its artefacts, and its documentation), how much of the union standard does it satisfy, and which regime does it fall short of?

\subsection{Scoring model}\label{sec:scoring}

Each of the twelve requirements of \Cref{tab:crosswalk} is scored on a pipeline as \emph{met}, \emph{partial}, or \emph{absent}. A requirement is \emph{met} only when the pipeline carries the artefact the strictest regime in that row demands, so for record retention (row 7), a pipeline is scored against the EU AI Act's $\geq$6-month floor even when deployed where no floor exists. From the twelve per-requirement scores we report two aggregates:

\begin{itemize}[leftmargin=*]
    \item a \textbf{per-regime score}: the fraction of that regime's applicable requirements the pipeline meets, so a firm can read its standing under SR 11-7, the EU AI Act, PRA SS1/23, MAS FEAT, or OSFI E-23 individually; and
    \item a \textbf{union score}: the fraction of the full twelve met at the strictest level, the single number that certifies a pipeline as deployable across all five jurisdictions without per-market rework.
\end{itemize}

Six of the twelve requirements are mechanically checkable from a repository using the same static-analysis approach validated in the predecessor rubric of \citet{KhanMrAudit2026}: inventory metadata (row 1), documentation and environment (row 2), data references (row 3), version control (row 4), the presence of a retention policy (row 7), and a decommissioning record (row 12). The remaining six require either code execution or manual review against a published codebook; the scoring harness automates the first set and structures the second.

\subsection{What is released}\label{sec:released}

The artefact accompanying this paper is released under permissive licences (CC-BY-4.0 for the crosswalk, MIT for code) at the companion repository and comprises: (i) the crosswalk of \Cref{tab:crosswalk} as a machine-readable table (\texttt{crosswalk.csv}), keyed by requirement and regime with the governing provision in each cell; (ii) a reference scoring rubric that operationalises each requirement as a checkable predicate; and (iii) a reference scoring harness with a public leaderboard, seeded on public-data machine-learning-finance pipelines so that the first entries are reproducible by any reader.

We are explicit about the version state: this is a v0.1 seed. The crosswalk is complete for the seven regimes studied; the harness automates the six mechanically checkable requirements and ships the codebook for the manual six; and the seed leaderboard is intentionally small and community-extensible rather than a comprehensive ranking. The benchmark is designed to be \emph{extended} (by adding regimes as columns, pipelines as rows, and updated provisions as the law changes), not to be a finished league table.

\subsection{Why a benchmark, not just a survey}\label{sec:why_benchmark}

A static survey of compliance goes out of date the moment a regulation changes. A versioned, open benchmark stays current: when OSFI E-23 takes effect in 2027, or when the U.S.\ revised guidance settles, a column is updated and every pipeline is re-scored without re-reading five texts. The benchmark also makes compliance \emph{comparable} across pipelines, which a prose survey cannot, and gives practitioners something to engineer against instead of a document to interpret.

\section{Discussion}\label{sec:discussion}

\subsection{How to use the crosswalk}\label{sec:use}

For a practitioner, \Cref{tab:crosswalk} is read column-up: identify the jurisdictions a model will touch, take the union of their non-empty cells, and that union is the requirement set to engineer against. The divergent rows (7, 11, 12) are the ones a single-regime team is most likely to miss. For a researcher, the crosswalk is a common rubric: a corpus of published pipelines can be scored against the union and against each regime, making ``would this work survive an audit?'' a measurable, comparable question rather than a rhetorical one.

\subsection{Limitations}\label{sec:limitations}

Three limitations bound the contribution honestly.

\begin{itemize}[leftmargin=*]
    \item \textbf{It is an interpretation, not legal advice.} Each cell records the most on-point public provision in our reading; a supervisor or counsel may map a requirement differently, and a ``---'' marks the absence of an \emph{explicit} analogue, not a determination that the regime permits the omission. Firms remain responsible for their own compliance conclusions.
    \item \textbf{Regulations move.} SR 11-7 / OCC 2011-12 was rescinded and replaced by revised interagency guidance in 2026; the EU AI Act's high-risk duties take effect 2 August 2026; OSFI E-23 takes effect 1 May 2027. The crosswalk is versioned to the texts as dated in \Cref{sec:crosswalk}, and the released artefact is built to be updated. A crosswalk that is not maintained is wrong within a year; maintenance is part of the contribution, not an afterthought.
    \item \textbf{The benchmark seed is small.} The v0.1 harness automates six of twelve requirements and the seed leaderboard is deliberately limited. The value at v0.1 is the rubric and the reproducible scoring method, not a comprehensive ranking; breadth is the community-extension path.
\end{itemize}

\subsection{Future work}\label{sec:future}

Three extensions are planned and are structured into the released artefact: adding further regimes as columns (e.g.\ Australia's APRA CPS 230, Hong Kong's HKMA guidance); automating more of the six manual requirements through code execution; and growing the leaderboard to a representative cross-section of public machine-learning-finance pipelines. Each is additive and leaves the existing crosswalk and scores stable.

\subsection{Conclusion}\label{sec:conclusion}

Machine-learning model-risk regulation in finance is converging on a common backbone --- inventory, replicable documentation, independent validation, monitoring, and named governance --- while diverging on retention, fairness, and decommissioning. We have made that structure explicit in a single crosswalk of twelve requirements across five binding regimes and two voluntary frameworks, and turned it into an open, versioned benchmark that scores a pipeline against the union of five jurisdictions at once. The crosswalk is the durable contribution; the benchmark is the instrument that keeps it useful as the law evolves.


\bibliographystyle{apalike}
\bibliography{bib/references}

\end{document}
